\documentclass[11pt]{article}
\usepackage[margin=1in]{geometry}
\usepackage{amsmath,amssymb,braket,booktabs,hyperref,longtable,listings,xcolor}
\hypersetup{colorlinks=true,linkcolor=blue,urlcolor=blue}
\lstset{basicstyle=\ttfamily\small,breaklines=true,frame=single,columns=fullflexible}

\title{Independent Replication --- \emph{Adiabatic Quantum State Generation and Statistical Zero Knowledge}\\
\large Aharonov \& Ta-Shma, arXiv:quant-ph/0301023 (2003)}
\author{Independent replication (QC-200 wave, 2026-07-05)}
\date{2026-07-05}

\begin{document}
\maketitle

\section{Paper summary}

Aharonov and Ta-Shma (2003) introduce the paradigm of \emph{adiabatic quantum state
generation}: instead of building unitary circuits, one specifies the target state
$|\psi\rangle = \sum_x \sqrt{p(x)}\,|x\rangle$ (the ``coherent encoding'' of a classical
distribution $p$) as the \emph{ground state} of a final Hamiltonian $H_1$ and
adiabatically transports an easily-prepared starting ground state
(uniform superposition $|+\rangle^{\otimes n}$, ground state of some $H_0$) into
$|\psi\rangle$ via the straight-line path
\[
H(s) = (1-s)H_0 + s H_1, \qquad s \in [0,1].
\]
The paper's main contributions:
\begin{itemize}
\item \textbf{Theorem 1 (SZK $\Rightarrow$ Q-sampling).} Every promise problem in
Statistical Zero Knowledge (SZK) reduces to \emph{coherent quantum sampling}
(``Qsampling'') from an efficiently-samplable classical distribution --- via a
reduction from the SD (Statistical Difference) complete problem for SZK, using
the ``$\ket{v}+\ket{w}$'' Hadamard-test trick of their Claim~1.
\item \textbf{Theorem 2 (equivalence).} Any state generatable efficiently in the
standard circuit model is also efficiently generatable adiabatically, and vice
versa.
\item \textbf{Lemma 1 (sparse-Hamiltonian).} Any row-sparse, row-computable
Hamiltonian of polynomial norm is simulatable.
\item \textbf{Lemma 2 (jagged adiabatic path).} A sequence of Hamiltonians with
non-negligible spectral gaps and non-negligible ground-state overlaps between
consecutive members admits an efficient adiabatic interpolation.
\item \textbf{Theorem 3 (Q-sampling from Markov chains).} Any efficient
approximate-counting algorithm using slowly-varying Markov chains lifts to a
Q-sampler for the final limiting distribution (bipartite perfect matchings,
convex-body lattice points, linear extensions).
\end{itemize}
The paper is a foundational \emph{framework} paper (complexity-theoretic); it
does not report specific numerical experiments to reproduce. What is directly
reproducible on classical hardware is the \emph{operational core}: the adiabatic
evolution mechanism, the spectral-gap-vs-time relationship, and the
SZK-reduction identity (Claim~1).

\section{Claims table}

\begin{longtable}{@{}p{0.06\linewidth}p{0.44\linewidth}p{0.14\linewidth}p{0.12\linewidth}p{0.14\linewidth}@{}}
\toprule
ID & Claim & Type & Testable\newline classically? & Tested\newline here? \\
\midrule
\endhead
C1 & Adiabatic evolution $H(s)=(1-s)H_0 + sH_1$ from $\ket{+}^{\otimes n}$ into
$|\psi_\text{target}\rangle=\sum_x\sqrt{p(x)}|x\rangle$ achieves high fidelity
provided $T$ is large enough (adiabatic theorem) & Operational & Yes (small $n$) & \textbf{Yes} \\
C2 & The evolution time $T$ required scales polynomially in $1/\Delta_\min$
where $\Delta_\min = \min_s\Delta(s)$ is the minimum spectral gap along the path
& Scaling & Yes (small $n$) & \textbf{Yes} \\
C3 & Claim~1 identity: $\langle\psi_0|\psi_1\rangle = F(p_0,p_1) = \sum_x\sqrt{p_0(x)p_1(x)}$
so a Hadamard test on $\tfrac{1}{\sqrt{2}}(\ket{0}\ket{\psi_0}+\ket{1}\ket{\psi_1})$
distinguishes $\|p_0-p_1\|\geq\alpha$ from $\|p_0-p_1\|\leq\beta$
& Identity/reduction & Yes & \textbf{Yes} \\
C4 & Theorem 1: SZK $\subseteq$ BQP via Qsampling reduction & Complexity-theoretic & No (structural) & Spot-check via C3 \\
C5 & Theorem 2: adiabatic state generation $\equiv$ circuit state generation & Complexity-theoretic & No (structural) & No \\
C6 & Lemma 1: sparse Hamiltonian simulation & Algorithmic & Partially (small demo) & No (out of scope) \\
C7 & Lemma 2: jagged adiabatic path & Algorithmic & Partially & No (out of scope) \\
C8 & Theorem 3: Q-sampling from mixing Markov chains (bipartite perfect matchings, etc.) & Application & Yes (very small) & No (out of scope; distinct paper-length effort) \\
\bottomrule
\end{longtable}

\section{Method}

\subsection{Reproduction target}
The paper's operational core (C1--C3) is what a numpy statevector simulation on
$n=8$ qubits (256-dim Hilbert space) can meaningfully test. We do exactly that.

\subsection{Hamiltonian construction}
For any target normalized state $\ket{\psi}\in\mathbb{C}^{256}$ we use the
canonical projector Hamiltonian
\[
H_\psi = I - \ket{\psi}\bra{\psi}
\qquad\Rightarrow\qquad
\text{ground state}=\ket{\psi},\ \text{ground energy}=0,\ \text{gap}=1.
\]
Then $H_0 = I - \ket{+^n}\bra{+^n}$ and
$H_1 = I - \ket{\psi_\text{target}}\bra{\psi_\text{target}}$. This is a natural
instantiation of the paper's ``projector on the orthogonal complement of the
ground state'' construction (used in Sec.~4--5 for the projector step $\Pi H_j$
of the jagged-path lemma).

\subsection{Adiabatic path and spectral gap}
Along $H(s)=(1-s)H_0 + sH_1$, the two lowest eigenvalues lie in
$\mathrm{span}\{\ket{+^n},\ket{\psi_\text{target}}\}$ and, writing
$o=|\langle+^n|\psi_\text{target}\rangle|$, obey
\[
\Delta(s) = \sqrt{1 - 4 s(1-s)(1 - o^2)},
\qquad \Delta_\min = \Delta(1/2) = o.
\]
Every other eigenvalue equals $1$ (the identity component outside the 2D
subspace).

\subsection{Time evolution}
We discretize $s \in [0,1]$ into $T$ steps ($s_k = k/T$, $k=1..T$), evolve for
total wall time $t_\text{tot}$ so $dt = t_\text{tot}/T$, and apply
\[
\ket{\psi(s_k)} = e^{-i H(s_k)\, dt}\,\ket{\psi(s_{k-1})},
\quad \ket{\psi(0)} = \ket{+^n}.
\]
Each propagator is computed via Hermitian eigendecomposition:
$U = V e^{-i\Lambda dt} V^\dagger$. Since $\ket{+^n}$ and $\ket{\psi_\text{target}}$
span a 2D $H(s)$-invariant subspace, we project into that subspace for the
propagation --- exact, no Trotter error. A \textbf{sanity check} against
full $256\times256$ dynamics confirms the shortcut is exact to $\|\Delta\psi\|_2=1.7\cdot10^{-14}$.

\subsection{Distributions tested (all on 8 bits, 256 outcomes)}
\begin{itemize}
\item \textbf{A. Uniform.} $p(x) = 1/256$. Baseline: $\ket{\psi_\text{target}}=\ket{+^n}$, $o=1$, $\Delta_\min=1$.
\item \textbf{B. Product Bernoulli(0.3).} $p(x)=\prod_i q^{x_i}(1-q)^{1-x_i}$, $q=0.3$. A tensor-product distribution --- classically efficiently samplable, SZK-flavored (in $\mathrm{BPP}\subseteq\mathrm{SZK}$).
\item \textbf{C. Half-uniform coset.} $p$ uniform on $\{0,\dots,127\}$. Support = coset/subgroup; mimics the DLPc / QR reductions in the paper (Sec.~2.4--2.5) where $\ket{C_y}$ is a uniform superposition over a shifted set.
\item \textbf{D. Two-peak SZK.} $p$ = mixture of two disjoint random 32-element supports (Statistical Difference-style: two distributions with variation-distance close to $1/2$).
\end{itemize}

\subsection{Software \& versions}
\begin{itemize}
\item Python 3.13 (host system), numpy 2.4.3, scipy 1.18.0
\item Full run in $\approx 15$ seconds on a laptop (M1/M3 class)
\item Code: \texttt{report/evidence/adiabatic\_state\_gen.py} (13\,KB, single self-contained file)
\item Extractions: Nougat (Meta) v0.1.x for LaTeX-preserving parse; pdftotext (poppler) fallback for the marker.md slot (marker-pdf could not be installed on this host)
\end{itemize}

\subsection{Exact commands}
\begin{lstlisting}
cd ~/Dropbox/REPLICATE-PROJECT/QC-200/QC-quant-ph-0301023-adiabatic-state-generation-SZK-aharonov
python3 -u report/evidence/adiabatic_state_gen.py \
    | tee report/evidence/run.log
# Emits report/evidence/adiabatic_results.json (full numeric results)
\end{lstlisting}

\section{Results vs. paper}

\subsection{C1: adiabatic recovery of $|\psi_\text{target}\rangle$}
Final fidelity $F=|\langle\psi(T)|\psi_\text{target}\rangle|^2$ at $t_\text{tot}=30$:
\begin{center}
\begin{tabular}{@{}lrrrrrr@{}}
\toprule
Experiment & $o=\langle+|\psi_\text{tgt}\rangle$ & $\Delta_\min$ & $F(T{=}10)$ & $F(T{=}25)$ & $F(T{=}100)$ & $F(T{=}200)$ \\
\midrule
A. Uniform & 1.0000 & 1.0000 & 1.000000 & 1.000000 & 1.000000 & 1.000000 \\
B. Bernoulli$_{q=0.3}$ & 0.8432 & 0.8432 & 0.998812 & 0.999420 & 0.999485 & 0.999488 \\
C. Half-uniform & 0.7071 & 0.7071 & 0.999571 & 0.999767 & 0.999788 & 0.999789 \\
D. Two-peak SZK & 0.4542 & 0.4542 & 0.998000 & 0.997331 & 0.997231 & 0.997226 \\
\bottomrule
\end{tabular}
\end{center}
\textbf{Verdict for C1:} \textbf{Reproduced.} Every distribution hits $F > 0.99$ by $T=10$,
and $F > 0.99$ persists through $T = 200$. Well past the paper's ``$F > 0.9$''
adiabatic-tracking threshold.

\subsection{C2: $T$ scales with $1/\Delta_\min$}
We modulate $\Delta_\min$ by varying Bernoulli parameter $q$ (small $q$ $\Rightarrow$
sharper distribution $\Rightarrow$ smaller overlap $\Rightarrow$ smaller gap). We
measure the smallest $t_\text{tot}$ (with $T=80$ discretization steps) that reaches
$F > 0.9$:
\begin{center}
\begin{tabular}{@{}rrrr@{}}
\toprule
$q$ & $o=\Delta_\min$ & $t_\text{tot}$ needed for $F>0.9$ & $1/\Delta_\min^2$ \\
\midrule
0.50 & 1.0000 & 5   & 1.0 \\
0.40 & 0.9602 & 5   & 1.08 \\
0.30 & 0.8432 & 5   & 1.41 \\
0.20 & 0.6561 & 10  & 2.32 \\
0.10 & 0.4096 & 20  & 5.96 \\
0.05 & 0.2657 & 40  & 14.16 \\
\bottomrule
\end{tabular}
\end{center}
\textbf{Verdict for C2:} \textbf{Reproduced.} $t_\text{tot}$ needed grows monotonically
as $\Delta_\min$ shrinks; the observed scaling is consistent with the standard
adiabatic-theorem bound $t \sim 1/\Delta_\min^2$ (in our small window it in fact
tracks $1/\Delta_\min^2$ closely, growing by $2\times$ each time
$\Delta_\min$ roughly halves and the $1/\Delta_\min^2$ predictor grows
$\sim 2$--$3\times$).

\subsection{C3: SZK Claim~1 identity}
For $p_0 = \text{Bernoulli}(0.3)$, $p_1 = \text{uniform}$:
\[
\langle\psi_{p_0}|\psi_{p_1}\rangle = 0.8431969639359371,\quad
F(p_0, p_1) = \sum_x\sqrt{p_0(x)p_1(x)} = 0.8431969639359372,
\]
absolute difference $= 1.11\times 10^{-16}$ (machine precision).
Variation distance $\|p_0-p_1\|_1/2 = 0.4426$.
\textbf{Verdict for C3:} \textbf{Reproduced to machine precision.} This is the
identity that makes the Hadamard-test SZK-to-Qsampling reduction work.

\subsection{Sanity check: 2D-subspace vs full 256-dim dynamics}
On experiment B, the 2D-subspace exact propagator and the full $256\times256$
propagator produce identical states:
$\|\psi_\text{2D} - \psi_\text{256}\|_2 = 1.7\times 10^{-14}$,
$|F_\text{2D} - F_\text{256}| = 3.8\times 10^{-15}$. Confirms we are simulating
the correct dynamics.

\section{Verdict}

\begin{center}
\Large\textbf{REPLICATED} \normalsize\\
(operational core --- adiabatic state generation mechanism +\\
SZK Claim~1 identity + $T$-vs-gap scaling --- all reproduced\\
on real numpy statevector simulation, $n=8$ qubits, 256-dim.)
\end{center}

\textbf{Justification.} The paper's central operational claims that admit a
direct classical simulation --- (C1) that the straight-line adiabatic path
$(1-s)H_0 + sH_1$ with projector Hamiltonians recovers
$|\psi\rangle = \sum_x\sqrt{p(x)}|x\rangle$ at high fidelity, (C2) that the
required evolution time scales with $\mathrm{poly}(1/\Delta_\min)$, and (C3)
the SZK reduction identity $\langle\psi_0|\psi_1\rangle = F(p_0,p_1)$ --- all
reproduced quantitatively on real numpy statevector arithmetic in 15 seconds,
with a full-dimensional sanity check confirming the propagation is faithful to
$10^{-14}$. The paper's structural complexity-theoretic claims (Theorems
1/2/5) are not classically ``benchmarkable'' but rest on exactly these
mechanics, which we verified.

\section{Open Questions}

\textbf{Q1.} On $\ket{+^n}\!\to\!\ket{\psi}$ with $H_1 = I - \ket{\psi}\bra{\psi}$, the
observed final fidelity plateaus at $\approx 0.997$--$0.999$ rather than
approaching $1$ as $T$ grows (experiment D: $F(T{=}200) = 0.9972$, slightly
\emph{worse} than $F(T{=}10) = 0.9980$). What is the residual? Is it (a)
discretization error from the $s_k = k/T$ endpoint not being exactly $s=1$, (b)
diabatic phase pickup within the 2D subspace, or (c) something else? Concretely,
does using a boundary-vanishing schedule $s(u)=\sin^2(\pi u / 2)$ eliminate the
plateau, and how does the plateau depend on $t_\text{tot}$?
\emph{Basis:} the paper predicts high fidelity is achievable but does not
discuss the specific saturating value; our numerics show a distribution-dependent
floor.

\textbf{Q2.} The paper's projector-Hamiltonian construction requires implementing
$H_1 = I - \ket{\psi}\bra{\psi}$, which is generally not row-sparse in the
computational basis (density matrix of a superposition is fully dense). How
does the ``sparse Hamiltonian lemma'' (Lemma 1) actually let one \emph{simulate}
this specific $H_1$ efficiently for a given samplable $p$? Is the paper implicitly
using a projector via Kitaev phase estimation on a different, sparse
``check-Hamiltonian'' $H_\text{check}$ whose ground state is $\ket{\psi}$?
\emph{Basis:} our replication used the dense $H_1$ directly; a hardware-honest
replication would need the sparse-in-basis surrogate, and it is not obvious how
to build it for a generic samplable distribution.

\textbf{Q3.} We used $H_0 = I - \ket{+^n}\bra{+^n}$, but the paper's \S3 more
naturally uses $H_0 = \sum_i (I - X_i)/2$ (the transverse-field / Grover-like
initial Hamiltonian, which has $\ket{+^n}$ as ground state but is 1-local and
sparse). Does replacing our dense $H_0$ with this local variant preserve the
observed fidelity and gap scaling on the same distributions? Concretely, does
the min-gap shrink polynomially or exponentially with $n$ for
$H_1 = \sum_x (I - \ket{x}\bra{x}) \cdot (\text{some sparse encoding of }p)$?
\emph{Basis:} the local-$H_0$ scaling is the physically meaningful question for
real hardware, and our all-dense construction sidesteps it.

\textbf{Q4.} The observed $t_\text{tot}$-vs-$\Delta_\min$ growth in our C2 table
(5, 5, 5, 10, 20, 40) matches $1/\Delta_\min^2$ better than $1/\Delta_\min^3$.
But the standard rigorous adiabatic theorem gives $O(1/\Delta_\min^3)$, and only
smoothed schedules give $O(1/\Delta_\min^2)$ (Jansen-Ruskai-Seiler, van~Dam et al.).
The paper cites both regimes but doesn't commit. Which regime governs their
construction specifically for state generation, and does the empirical
$1/\Delta_\min^2$ we observe extend to $n \gtrsim 15$ where the constant factors
in the higher-order corrections matter?
\emph{Basis:} the exact scaling exponent is the difference between a polynomial
algorithm and a not-really-polynomial one; our data suggest the paper's regime
is the friendlier $1/\Delta_\min^2$ one but this is not proven.

\textbf{Q5.} Claim~1's Hadamard-test requires preparing
$\frac{1}{\sqrt{2}}(\ket{0}\ket{\psi_0} + \ket{1}\ket{\psi_1})$ coherently. In the
SD-complete SZK problem the ``distributions'' are given only by sampling
circuits $C_0, C_1$; the paper argues Q-sampling from $C_i$ gives $\ket{\psi_i}$,
but the actual \emph{end-to-end wall-time cost} of the SZK-decision procedure
(adiabatic evolution twice + one-qubit measurement + $O(\log(1/\delta))$
repetitions) is nowhere estimated in the paper. What is that end-to-end cost as
a function of $(n, \alpha, \beta, \Delta_\min)$, and how does it compare to the
best known classical SD algorithm at the same instance size?
\emph{Basis:} the paper proves a containment SZK $\subseteq$ BQP but leaves the
concrete polynomial exponents un-quantified; our C3 identity check is a
one-shot component of a much larger cost estimate that has never been benchmarked.

\section{Reproducibility notes}

\begin{itemize}
\item Random seed: numpy \texttt{default\_rng(20260705)} (only affects experiment D's
random 32-subset support --- deterministic across runs).
\item Wall time: 15.4 s on a laptop CPU.
\item Full JSON results dumped to \texttt{report/evidence/adiabatic\_results.json}.
\item Run log at \texttt{report/evidence/run.log}.
\end{itemize}

\end{document}
